% Compile with: uplatex file.tex && dvipdfmx file.dvi
%           or: latexmk -uplatex file.tex
\documentclass[dvipdfmx]{article}
\usepackage{amsmath, amssymb}
\usepackage{tikz}

\title{scenario\_damped\_oscillator\_numeric}
\date{\today}

\begin{document}
\maketitle

First, prepare the time-dependent displacement $x$ and the damping rate $\gamma$ / natural angular frequency $\omega_0$ (both positive reals).

Define x as a time-dependent function

$\checkmark$ Declared x as time-dependent (time variable: t)

Define $\gamma$ as a positive real number

\[
  \gamma
\]

\noindent\textbf{Assumption:}\ $\displaystyle \gamma > 0$\par

Define $\omega_0$ as a positive real number

\[
  \omega_{0}
\]

\noindent\textbf{Assumption:}\ $\displaystyle \omega_{0} > 0$\par

Consider the damped oscillator's equation of motion below ($\gamma$ is the damping rate, $\omega_0$ is the natural angular frequency).

$\ddot{x} + 2 \gamma \dot{x} + \omega_0^2 x = 0$

\begin{equation}
  \ddot{x} + 2 \gamma \dot{x} + \omega_0^2 x = 0
  \label{eq:E1}
\end{equation}

Substituting the exponential ansatz $x=e^{rt}$ into this differential equation ($\dot x=re^{rt},\ \ddot x=r^2e^{rt}$, then dividing both sides by $e^{rt}\neq0$) gives the characteristic equation of this differential equation below.

$r^2 + 2 \gamma r + \omega_0^2 = 0$

\begin{equation}
  r^2 + 2 \gamma r + \omega_0^2 = 0
  \label{eq:E2}
\end{equation}

Since the characteristic equation is a quadratic in $r$ ($ar^2+br+c=0$ with $a=1,\ b=2\gamma,\ c=\omega_0^2$), its discriminant is $b^2-4ac=(2\gamma)^2-4\omega_0^2$. To phrase the discussion below in the natural form "positive means underdamped", we define the sign-flipped quantity $D=4\omega_0^2-(2\gamma)^2=-(b^2-4ac)$.

$D = 4 \omega_0^2 - (2 \gamma)^2$

\begin{equation}
  D = 4 \omega_0^2 - (2 \gamma)^2
  \label{eq:E3}
\end{equation}

Here we assume underdamped (oscillatory) motion, $D>0$. This choice is a physical judgment that depends on the problem being treated, and is not something the tool should decide on its own, so it is assumed explicitly, as follows.

Assume D $>$ 0

\noindent\textbf{Assumption:}\ $\displaystyle D > 0$\par

Under the assumption $D>0$, find the general solution of the equation of motion (E1).

Find the general solution of E1

\begin{equation}
  x=e^ {- \left(\gamma t\right) } \left(k_{1} \sin \left({{\sqrt{4 \omega_{0}^2-4 \gamma^2} t}\over{2}}\right)+k_{2} \cos \left({{\sqrt{4 \omega_{0}^2-4 \gamma^2} t}\over{2}}\right)\right)
  \label{eq:E4}
\end{equation}

The general solution (E4) contains arbitrary constants $k_1,k_2$. Give the initial conditions $x(0)=1,\ \dot x(0)=0$ to fix them.

Set the initial conditions of E4 to x(0)=1, x'(0)=0

\begin{equation}
  x = e^ {- \left(\gamma t\right) } \left(\cos \left({{\sqrt{4 \omega_{0}^2-4 \gamma^2} t}\over{2}}\right)-{{\gamma \sqrt{4 \omega_{0}^2-4 \gamma^2} \sin \left({{\sqrt{4 \omega_{0}^2-4 \gamma^2} t}\over{2}}\right)}\over{2 \gamma^2-2 \omega_{0}^2}}\right)
  \label{eq:E5}
\end{equation}

To fork into the numeric layer in the next step, extract and simplify just the right-hand side of the determined solution (E5).

Simplify the right-hand side of E5

\begin{equation}
  x = -{{\gamma \sqrt{4 \omega_{0}^2-4 \gamma^2} \sin \left({{\sqrt{4 \omega_{0}^2-4 \gamma^2} t}\over{2}}\right)+\left(2 \omega_{0}^2-2 \gamma^2\right) \cos \left({{\sqrt{4 \omega_{0}^2-4 \gamma^2} t}\over{2}}\right)}\over{2 e^{\gamma t} \gamma^2-2 e^{\gamma t} \omega_{0}^2}}
  \label{eq:E6}
\end{equation}

\end{document}